\section{Môi trường và Công nghệ Triển khai}

\subsection{Cấu hình Môi trường Hệ thống}

\begin{enumerate}
    \item Cơ sở lý thuyết về Tương thích Nền tảng (Platform Compatibility) và I/O Phi đồng bộ
    
    Trong kiến trúc phần mềm mạng, cơ chế Vào/Ra (I/O) là yếu tố quan trọng đối với năng lực xử lý đồng thời. Hệ điều hành nhân Linux/Unix cung cấp API \texttt{epoll} với độ phức tạp $O(1)$, cải thiện hiệu suất trong việc quản lý số lượng lớn kết nối (Socket) so với hàm \texttt{select()} có độ phức tạp $O(N)$. Nền tảng Windows sử dụng kiến trúc \texttt{IOCP} (I/O Completion Ports). Do sự khác biệt về API giữa các hệ điều hành, phân hệ máy chủ (Server) được thiết kế cho môi trường Linux nhằm tối ưu hóa tốc độ xử lý.
    
    Quy trình đóng gói (Build Process) sử dụng Hệ thống Siêu biên dịch (Meta-Build System) CMake nhằm chuẩn hóa quy trình triển khai. CMake cung cấp lớp trừu tượng (Abstraction Layer) cho phép tự động khởi tạo kịch bản biên dịch bản địa, thay thế cho việc duy trì các tệp \texttt{Makefile} độc lập trên từng nền tảng.

    \item Kiến trúc Biên dịch trong tệp \texttt{CMakeLists.txt}
    
    Dự án được tổ chức theo cấu trúc Kho lưu trữ tập trung (Monorepo), chứa toàn bộ mã nguồn của \texttt{vcs\_server} và \texttt{vcs\_client}. Tiến trình biên dịch được cấu hình thông qua các tham số của CMake:
    
    \begin{itemize}[label={-}]
        \item Phân hệ Máy chủ (Server-Side): Triển khai trên môi trường POSIX (Linux). Cấu trúc kiểm tra \texttt{if(NOT (WIN32 OR CMAKE\_CROSSCOMPILING))} được sử dụng để thiết lập tính tương thích và đảm bảo khả năng hoạt động của module \texttt{epoll}.
        
        \item Phân hệ Máy khách (Client-Side): Hỗ trợ kiến trúc đa nền tảng. Khi biến \texttt{WIN32} được kích hoạt, CMake thực thi liên kết tĩnh (Static Linking) với các thư viện mạng cốt lõi của Windows như \texttt{wsock32}, \texttt{ws2\_32} và \texttt{crypt32}.
        
        \item Quản lý Phụ thuộc (Dependency Management): Dự án định cấu hình biên dịch dựa trên chuẩn \texttt{C++17} (\texttt{CMAKE\_CXX\_STANDARD 17}) kết hợp cùng các cờ tối ưu hóa \texttt{-Wall -Wextra -O2}. Quá trình liên kết nhị phân phụ thuộc vào 4 thư viện lõi: OpenSSL (\texttt{libssl-dev}), SQLite3 (\texttt{libsqlite3-dev}), POSIX Threads (Threads) và \texttt{nlohmann/json}.
    \end{itemize}
\end{enumerate}

Bảng \ref{tab:dependencies} trình bày ma trận các thành phần phụ thuộc (Dependencies) trong kiến trúc phần mềm.

\begin{longtable}{
    >{\raggedright\arraybackslash}p{4cm}
    >{\centering\arraybackslash}p{3cm}
    >{\raggedright\arraybackslash}p{7cm}
}
    \caption[Bảng Ma trận Phụ thuộc Thư viện]{Bảng Ma trận Phụ thuộc Thư viện (Dependency Matrix)} \label{tab:dependencies} \\
    \toprule
    \multicolumn{1}{c}{\textbf{Tên Thư viện}} & \multicolumn{1}{c}{\textbf{Loại Liên kết}} & \multicolumn{1}{c}{\textbf{Mục đích Kỹ thuật}} \\
    \midrule
    \endfirsthead
    
    \toprule
    \multicolumn{1}{c}{\textbf{Tên Thư viện}} & \multicolumn{1}{c}{\textbf{Loại Liên kết}} & \multicolumn{1}{c}{\textbf{Mục đích Kỹ thuật}} \\
    \midrule
    \endhead
    
    \bottomrule
    \endfoot
    
    \bottomrule
    \endlastfoot

    \texttt{OpenSSL::SSL} \newline \texttt{OpenSSL::Crypto} & Động (Dynamic) & Cung cấp phương thức mã hóa lai (RSA, AES-GCM) và hàm sinh số ngẫu nhiên CSPRNG. \\

    \texttt{SQLite3} & Động (Dynamic) & Hệ quản trị cơ sở dữ liệu cục bộ phục vụ phân hệ \texttt{AuthManager} và \texttt{ChatHistory}. \\

    \texttt{Threads::Threads} & Hệ thống (System) & Gọi API POSIX Threads nhằm thực thi đa luồng cho các Worker Threads và vòng lặp sự kiện. \\

    \texttt{nlohmann/json} & Header-only & Hỗ trợ tuần tự hóa (Serialize) và giải tuần tự hóa (Deserialize) chuỗi JSON. \\

    \texttt{wsock32}, \texttt{ws2\_32} & Tĩnh (Static) \newline (Chỉ Windows) & API mạng lõi thay thế cho Socket POSIX ở trình khách trên nền tảng hệ điều hành Windows. \\
\end{longtable}

\subsection{Kịch bản Biên dịch và Vận hành}

\subsubsection{Cơ sở lý thuyết về quy trình biên dịch}

Trong kỹ thuật phần mềm hiện đại, quy trình xây dựng (Build Process) hỗ trợ chuẩn hóa mã nguồn trên đa nền tảng. CMake là một Hệ thống Siêu biên dịch (Meta-build System), có chức năng chuyển đổi cấu hình khai báo (Declarative Configuration) thành các tập lệnh thực thi (Native Build Scripts) tương thích với môi trường biên dịch mục tiêu.

Dự án áp dụng cấu trúc Biên dịch ngoại tuyến (Out-of-Source Build). Các tệp đối tượng trung gian (\texttt{.o}, \texttt{.so}) được lưu trữ tại thư mục \texttt{build/}, giúp giữ nguyên cấu trúc mã nguồn gốc. Kỹ thuật Biên dịch song song (Parallel Compilation) thông qua bộ đa xử lý đối xứng (SMP) của CPU được sử dụng nhằm giảm thời gian đóng gói nhị phân.

\subsubsection{Kiến trúc biên dịch thực tế}

Tiến trình đóng gói được cấu hình qua tệp \texttt{CMakeLists.txt} và tệp kịch bản \texttt{build\_all.sh}. Hệ thống yêu cầu trình biên dịch chuẩn \texttt{C++17} kết hợp với các tham số \texttt{-Wall -Wextra -O2} nhằm kiểm soát các cảnh báo và tối ưu hóa hiệu năng thực thi.

Quá trình biên dịch phân tách thành hai luồng độc lập:
\begin{itemize}[label={-}]
    \item Phân luồng Máy chủ (Server-side): Yêu cầu hệ điều hành POSIX (Linux) thông qua khối kiểm tra \texttt{if(NOT (WIN32 OR CMAKE\_CROSSCOMPILING))}. Mã nhị phân được liên kết với POSIX Threads, OpenSSL và \texttt{SQLITE3}.
    \item Phân luồng Máy khách (Client-side): Hỗ trợ Biên dịch chéo (Cross-platform Compilation). Đối với kiến trúc Windows, tệp \texttt{mingw-toolchain.cmake} thực hiện chức năng cấu hình MinGW nhằm liên kết tĩnh các thư viện mạng đặc thù (\texttt{wsock32}, \texttt{ws2\_32}, \texttt{crypt32}).
\end{itemize}

Bảng \ref{tab:build_commands} chi tiết hóa các tập lệnh cấu hình trong quy trình triển khai.

\begin{longtable}{
    >{\raggedright\arraybackslash}p{4cm}
    >{\raggedright\arraybackslash}p{5cm}
    >{\raggedright\arraybackslash}p{5.5cm}
}
    \caption{Các tập lệnh và tham số cấu hình trong quy trình biên dịch} \label{tab:build_commands} \\
    \toprule
    \multicolumn{1}{c}{\textbf{Giai đoạn}} & \multicolumn{1}{c}{\textbf{Tập lệnh (Command)}} & \multicolumn{1}{c}{\textbf{Mục đích kỹ thuật}} \\
    \midrule
    \endfirsthead
    
    \toprule
    \multicolumn{1}{c}{\textbf{Giai đoạn}} & \multicolumn{1}{c}{\textbf{Tập lệnh (Command)}} & \multicolumn{1}{c}{\textbf{Mục đích kỹ thuật}} \\
    \midrule
    \endhead
    
    \bottomrule
    \endfoot
    
    \bottomrule
    \endlastfoot

    1. Khởi tạo Không gian & \texttt{mkdir -p build/linux \&\& cd build/linux} & Thiết lập môi trường cho phương thức Biên dịch ngoại tuyến (Out-of-source Build). \\

    2. Thiết lập CMake & \texttt{cmake ../..} & Khởi chạy trình phân tích, định tuyến thư viện và cấu trúc tệp \texttt{Makefile}. \\

    3. Biên dịch Đa luồng & \texttt{make -j\$(nproc)} & Kích hoạt tác vụ biên dịch song song dựa trên số lượng nhân CPU. \\

    4. Biên dịch Chéo (Win) & \texttt{cmake -DCMAKE\_TOOLCHAIN...} \newline \texttt{make -j\$(nproc) vcs\_client} & Khởi chạy kịch bản biên dịch ứng dụng Client độc lập cho hệ điều hành Windows. \\
\end{longtable}


\subsubsection{Chu trình vận hành và tham số dòng lệnh}

Các tập tin nhị phân hỗ trợ hệ thống phân tích tham số dòng lệnh (Command-line Arguments). Chức năng này cho phép cấu hình hệ thống tại thời điểm chạy (Runtime) mà không cần biên dịch lại mã nguồn.

1. Phân hệ Máy chủ (Server Execution)

Tiến trình \texttt{vcs\_server} thực hiện phân bổ Thread Pool, cấp phát cơ sở dữ liệu và mở cổng mạng chờ kết nối. Hàm khởi tạo \texttt{main()} thiết lập thư mục \texttt{logs/} (phân quyền \texttt{0755}) và nạp cơ sở dữ liệu \texttt{vcs\_chat.db}. Bảng \ref{tab:server_args} mô tả chức năng của các tham số cấu hình máy chủ.

\begin{longtable}{
    >{\centering\arraybackslash}p{4cm}
    >{\centering\arraybackslash}p{3cm}
    >{\raggedright\arraybackslash}p{7cm}
}
    \caption{Các tham số dòng lệnh vận hành Máy chủ} \label{tab:server_args} \\
    \toprule
    \multicolumn{1}{c}{\textbf{Tham số}} & \multicolumn{1}{c}{\textbf{Kiểu dữ liệu}} & \multicolumn{1}{c}{\textbf{Mô tả chức năng kỹ thuật}} \\
    \midrule
    \endfirsthead
    
    \toprule
    \multicolumn{1}{c}{\textbf{Tham số}} & \multicolumn{1}{c}{\textbf{Kiểu dữ liệu}} & \multicolumn{1}{c}{\textbf{Mô tả chức năng kỹ thuật}} \\
    \midrule
    \endhead
    
    \bottomrule
    \endfoot
    
    \bottomrule
    \endlastfoot

    \texttt{--config <path>} & Chuỗi & Chỉ định tệp cấu hình JSON chứa thông số vận hành (Mặc định: \texttt{server\_config.json}). \\

    \texttt{--port <port\_number>} & Số nguyên & Ghi đè tham số cổng dịch vụ mạng khai báo trong cấu hình hệ thống. \\

    \texttt{--verify-audit-log} & Cờ (Bool) & Kích hoạt phương thức \texttt{AuditLogger::verifyChain()} nhằm xác thực tính toàn vẹn của chuỗi Audit Log. \\
\end{longtable}

2. Phân hệ Máy khách (Client Execution)

Ứng dụng \texttt{vcs\_client} khởi tạo liên kết Socket thông qua cấu trúc \texttt{ConnectionManager}. Đối với nền tảng Windows, API \texttt{WSAStartup} được gọi ở giai đoạn đầu nhằm khởi động thư viện mạng \texttt{winsock2}. Bảng \ref{tab:client_args} liệt kê các tham số dòng lệnh liên quan.

\begin{longtable}{
    >{\centering\arraybackslash}p{4cm}
    >{\centering\arraybackslash}p{3cm}
    >{\raggedright\arraybackslash}p{7cm}
}
    \caption{Các tham số dòng lệnh Máy khách} \label{tab:client_args} \\
    \toprule
    \multicolumn{1}{c}{\textbf{Tham số}} & \multicolumn{1}{c}{\textbf{Kiểu dữ liệu}} & \multicolumn{1}{c}{\textbf{Mô tả chức năng kỹ thuật}} \\
    \midrule
    \endfirsthead
    
    \toprule
    \multicolumn{1}{c}{\textbf{Tham số}} & \multicolumn{1}{c}{\textbf{Kiểu dữ liệu}} & \multicolumn{1}{c}{\textbf{Mô tả chức năng kỹ thuật}} \\
    \midrule
    \endhead
    
    \bottomrule
    \endfoot
    
    \bottomrule
    \endlastfoot

    \texttt{--host <ip\_address>} & Chuỗi & Định tuyến đến địa chỉ IP máy chủ (Mặc định: \texttt{127.0.0.1}). \\

    \texttt{--port <port\_number>} & Số nguyên & Chỉ định cổng mạng của máy chủ đích (Mặc định: \texttt{9000}). \\
\end{longtable}

